\section{RNN 的梯度问题}

\subsection{梯度消失问题}

训练 RNN 时，需要将损失通过时间反向传播（Backpropagation Through Time, BPTT）。考虑在位置 $t$ 处的损失 $J^{(t)}(\theta)$，其相对于早期隐藏状态 $h^{(k)}$ 的梯度涉及一连串偏导数的乘积：

$$\frac{\partial J^{(t)}}{\partial h^{(k)}} = \frac{\partial h^{(k+1)}}{\partial h^{(k)}} \times \frac{\partial h^{(k+2)}}{\partial h^{(k+1)}} \times \cdots \times \frac{\partial J^{(t)}}{\partial h^{(t)}}$$

\begin{figure}[H]
\centering
\includegraphics[width=0.95\textwidth]{slides-images/slide-005.jpg}
\caption{RNN 梯度消失与梯度爆炸问题：反向传播时梯度需要经过多次矩阵乘法\protect\footnotemark}
\end{figure}
\footnotetext{Slides 第6页。}

\begin{figure}[H]
\centering
\includegraphics[width=0.95\textwidth]{slides-images/slide-010.jpg}
\caption{梯度消失的直觉：当偏导数 $\frac{\partial h^{(k+1)}}{\partial h^{(k)}}$ 较小时，梯度信号随反向传播距离指数级衰减\protect\footnotemark}
\end{figure}
\footnotetext{Slides 第11页。}

\begin{importantbox}{梯度消失的数学根源}
假设忽略非线性激活函数，$h^{(t)} = W_h h^{(t-1)} + W_x x^{(t)}$，则：
$$\frac{\partial h^{(t)}}{\partial h^{(t-1)}} = W_h$$
反向传播 $n$ 步等于计算 $W_h^n$。对 $W_h$ 进行特征值分解：
\begin{itemize}
    \item 若所有特征值 $|\lambda_i| < 1$：$W_h^n \to 0$（梯度消失）
    \item 若存在特征值 $|\lambda_i| > 1$：$W_h^n \to \infty$（梯度爆炸）
\end{itemize}
\end{importantbox}

\subsection{梯度消失的实际影响}

实验表明，简单 RNN 的有效上下文窗口约为 \textbf{7 个 token}——更远的信息几乎无法被学习到。这与 5-gram 语言模型相比没有实质性进步。

\begin{figure}[H]
\centering
\includegraphics[width=0.95\textwidth]{slides-images/slide-012.jpg}
\caption{长距离依赖示例：预测``ticket''需要回溯 20 多个词，简单 RNN 无法捕捉这种依赖\protect\footnotemark}
\end{figure}
\footnotetext{Slides 第13页。}

考虑这样一段文本：``When she tried to print her ticket, she found that the printer was out of toner. She went to the store to buy more toner. ... After installing the toner into the printer, she finally printed her \underline{\hspace{1cm}}.'' 人类可以几乎 100\% 确定地预测最后一个词是``ticket''，但这依赖于 20 多个词之前的信息。简单 RNN 因梯度消失而无法建立这种长距离关联。

\subsection{梯度爆炸与梯度裁剪}

\begin{figure}[H]
\centering
\includegraphics[width=0.95\textwidth]{slides-images/slide-015.jpg}
\caption{梯度爆炸的危害：过大的梯度导致参数更新步长过大，模型可能``跑飞''到损失函数的糟糕区域\protect\footnotemark}
\end{figure}
\footnotetext{Slides 第16页。}

梯度爆炸的解决方案相对简单——\textbf{梯度裁剪}（Gradient Clipping）：

$$\hat{g} = \begin{cases} g & \text{if } \|g\| \leq \text{threshold} \\ \frac{\text{threshold}}{\|g\|} \cdot g & \text{if } \|g\| > \text{threshold} \end{cases}$$

\begin{figure}[H]
\centering
\includegraphics[width=0.95\textwidth]{slides-images/slide-017.jpg}
\caption{梯度裁剪：当梯度范数超过阈值时，按比例缩放梯度方向不变但大小受限\protect\footnotemark}
\end{figure}
\footnotetext{Slides 第18页。}

\begin{importantbox}{梯度裁剪是训练神经网络的必备技巧}
梯度裁剪虽然是一个``粗暴的 hack''，但它非常有效且几乎是必需的。阈值通常设为 5、10 或 20。梯度裁剪保持了梯度的方向，只限制其大小，避免了过大步长导致的训练不稳定。在训练任何深度模型时，都应该考虑使用梯度裁剪。
\end{importantbox}

\subsection{本章小结}

\begin{itemize}
    \item \textbf{梯度消失}：梯度信号在长序列中指数级衰减，导致模型无法学习长距离依赖（有效窗口仅 $\sim$7 个 token）
    \item \textbf{梯度爆炸}：梯度信号指数级增长，导致参数更新失控，可通过梯度裁剪解决
    \item 梯度消失是更根本的问题，需要从\textbf{架构层面}解决——这正是 LSTM 的动机
\end{itemize}

%% ========== LSTM ==========

